\documentclass[12pt]{article}
\usepackage[utf8]{inputenc}
\usepackage[T1]{fontenc}
\usepackage{amsmath,amssymb,amsthm}
\usepackage{geometry}
\geometry{margin=1in}
\usepackage{enumitem}
\setlist{nosep}

\linespread{1.1} % Slightly increased line spacing for better readability
% -------------------- Macros --------------------
\newcommand{\N}{\mathbb{N}}
\newcommand{\Z}{\mathbb{Z}}
\newcommand{\Q}{\mathbb{Q}}
\newcommand{\R}{\mathbb{R}}
\newcommand{\C}{\mathbb{C}}
\newcommand{\K}{\mathbb{K}} % used to denote either R or C
\newcommand{\eps}{\varepsilon}
\newcommand{\dd}{\,\mathrm{d}}

\newcommand{\NumberedDefinition}[3]{% #1 number, #2 title, #3 body
\paragraph*{Definition #1 ( #2 ).} #3\par}
\newcommand{\NumberedTheorem}[3]{% #1 number, #2 title, #3 body
\paragraph*{Theorem #1 ( #2 ).} #3\par}
\newcommand{\NumberedProposition}[3]{% #1 number, #2 title, #3 body
\paragraph*{Proposition #1 ( #2 ).} #3\par}
\newcommand{\NumberedLemma}[3]{% #1 number, #2 title, #3 body
\paragraph*{Lemma #1 ( #2 ).} #3\par}
\newcommand{\NumberedCorollary}[3]{% #1 number, #2 title, #3 body
\paragraph*{Corollary #1 ( #2 ).} #3\par}
\newcommand{\NumberedRemark}[3]{% #1 number, #2 title, #3 body
\paragraph*{Remark #1 ( #2 ).} #3\par}
\newcommand{\NumberedExample}[3]{% #1 number, #2 title, #3 body
\paragraph*{Example #1 ( #2 ).} #3\par}

% This chapter translation: Original German "Die Exponentialfunktion" (The Exponential Function)
% Numbering (Definition/Satz) preserved exactly as in source. Do NOT renumber.

\begin{document}
\subsection*{3.6. Power Series (Potenzreihen)}

\NumberedDefinition{3.6.1}{Power series}{Let $(a_n)_{n\in\N}$ be a sequence in $\K$ and $z_0\in\K$. The series
\[
	\sum a_n\,(z-z_0)^{n}
\]
in the variable $z\in\K$ is called a \emph{power series} with coefficients $a_n$ and center $z_0$.}

\NumberedExample{3.6.2}{}{
\begin{enumerate}[label={(\arabic*)}, leftmargin=*]
	\item $\displaystyle \sum \frac{z^{n}}{n!}$. This power series $(a_n=1/n!,\ z_0=0)$ converges absolutely for all $z\in\C$.
	\item $\displaystyle \sum z^{n}$. This power series $(a_n=1,\ z_0=0)$ converges absolutely for all $z\in\C$ with $|z|<1$.
	\item $\displaystyle \sum (n!)\, z^{n}$ with $a_n=n!,\ z_0=0$. Since
	\[
		\Bigg|\frac{(n+1)\,z^{n+1}}{n!\,z^{n}}\Bigg| = (n+1)|z| \longrightarrow \infty \quad (n\to\infty) \quad \text{for every } z\ne0,
	\]
	the series converges only for $z=0$.
\end{enumerate}}

\NumberedTheorem{3.6.3}{Radius of convergence}{Let $\sum a_n (z-z_0)^n$ be a power series and set
\[
	\rho := \frac{1}{\displaystyle \limsup_{n\to\infty} \sqrt[n]{|a_n|}}\;\in [0,\infty].
\]
Then:
\begin{enumerate}[label={(\arabic*)}, leftmargin=*]
	\item The power series converges absolutely for all $z\in\K$ with $|z-z_0|<\rho$.
	\item The power series diverges for all $z\in\K$ with $|z-z_0|>\rho$.
\end{enumerate}}
The number $\rho$ is called the \emph{radius of convergence} of the power series, and the set $\{z\in\K:\ |z-z_0|<\rho\}$ is called its (open) disc of convergence.

\paragraph{Proof.}
Use the root test:
\[
	\limsup_{n\to\infty} \sqrt[n]{\,|a_n (z-z_0)^n|\,}
	 = |z-z_0|\, \limsup_{n\to\infty} \sqrt[n]{|a_n|}
	 = \frac{|z-z_0|}{\rho}.
\]
Hence the series converges for $|z-z_0|<\rho$. The divergence for $|z-z_0|>\rho$ follows similarly. \qed

\NumberedTheorem{3.6.4}{Ratio form for radius}{Let $\sum a_n (z-z_0)^n$ be a power series and assume
\[
	\lim_{n\to\infty} \frac{|a_n|}{|a_{n+1}|} =: \alpha \in \overline{\R}.
\]
Then the radius of convergence satisfies $\rho = \alpha$.}
\paragraph{Proof.}
Apply the ratio test:
\[
	\Bigg|\frac{a_{n+1}(z-z_0)^{n+1}}{a_n (z-z_0)^n}\Bigg|
	 = |z-z_0|\, \Bigg|\frac{a_{n+1}}{a_n}\Bigg| \longrightarrow \frac{|z-z_0|}{\alpha}.
\]
Thus the series converges for $|z-z_0|<\alpha$ and does not converge for $|z-z_0|>\alpha$, so $\rho=\alpha$. \qed

\NumberedExample{3.6.5}{}{
\begin{enumerate}[label={(\arabic*)}, leftmargin=*]
	\item $\displaystyle \sum \frac{z^{n}}{n!}$: Here $\frac{|a_n|}{|a_{n+1}|} = \frac{(n+1)!}{n!} = n+1 \to \infty$, hence $\rho=\infty$.
	\item $\displaystyle \sum n^{m} z^{n}$ with fixed $m\in\Z$: Then
	\[
		\frac{|a_n|}{|a_{n+1}|} = \frac{n^{m}}{(n+1)^{m}} = \Big(\frac{n}{n+1}\Big)^{m} = \Big( \frac{1}{1+\tfrac{1}{n}} \Big)^{m} \longrightarrow 1,
	\]
	so $\rho=1$.
\end{enumerate}}

\NumberedRemark{3.6.6}{}{No general statement can be made about behavior on the boundary of the disc of convergence. For example:
\begin{enumerate}[label={(\arabic*)}, leftmargin=*]
	\item $\sum z^{n}$ has $\rho=1$ and diverges for every $z$ with $|z|=1$.
	\item $\sum \dfrac{z^{n}}{n}$ has $\rho=1$ and converges for $z=-1$ but diverges for $z=1$.
	\item $\sum \dfrac{z^{n}}{n^{2}}$ has $\rho=1$ and converges for all $|z|=1$ (convergent majorant $\sum \tfrac{1}{n^{2}}$).
\end{enumerate}}

\paragraph{Rules for power series.}

\NumberedTheorem{3.6.7}{Algebra of power series}{Let $\alpha\in\K$ and let $\sum a_n (z-z_0)^n$, $\sum b_n (z-z_0)^n$ be two power series with the same center $z_0$ and radii of convergence $\rho_1,\rho_2$. Then:
\begin{enumerate}[label={(\arabic*)}, leftmargin=*]
	\item $\displaystyle \sum_{n=0}^{\infty} \alpha a_n (z-z_0)^n = \alpha \sum_{n=0}^{\infty} a_n (z-z_0)^n$ for all $|z-z_0|<\rho_1$.
	\item $\displaystyle \sum_{n=0}^{\infty} (a_n+b_n)(z-z_0)^n = \sum_{n=0}^{\infty} a_n (z-z_0)^n + \sum_{n=0}^{\infty} b_n (z-z_0)^n$ for all $|z-z_0|<\min\{\rho_1,\rho_2\}$.
	\item $\displaystyle \sum_{n=0}^{\infty} \Big( \sum_{k=0}^{n} a_k b_{n-k} \Big) (z-z_0)^n 
		= \Big( \sum_{n=0}^{\infty} a_n (z-z_0)^n \Big) \Big( \sum_{n=0}^{\infty} b_n (z-z_0)^n \Big)$ for all $|z-z_0|<\min\{\rho_1,\rho_2\}$.
\end{enumerate}}

\paragraph{Function vs. series representation.}
Consider $f(z) := \dfrac{1}{1-z}$ for $z\in\C\setminus\{1\}$. For $|z|<1$ the function also has the representation as a power series
\[
	g(z) := \sum z^{n}.
\]
For $|z|\ge1$ this representation is no longer valid. However, for $|z|>1$, using
\[
	f(z) = -\frac{1}{z}\,\frac{1}{1-\tfrac{1}{z}} \qquad\text{and}\qquad \Big|\frac{1}{z}\Big|<1,
\]
we obtain the representation
\[
	f(z) = -\frac{1}{z}\sum_{i=0}^{\infty} \Big(\frac{1}{z}\Big)^{i} = -\sum_{i=0}^{\infty} z^{-i-1}.
\]


\end{document}